\documentclass[11pt,a4paper]{article}
\usepackage{fontspec}
\setmainfont{Noto Serif}
\setsansfont{Noto Sans}
\usepackage{babel}
\babelprovide[main,import]{russian}
\usepackage{amsmath,amssymb,mathtools,booktabs,geometry,fancyhdr,microtype,xcolor,hyperref}
\geometry{left=21mm,right=21mm,top=21mm,bottom=21mm}
\definecolor{ink}{HTML}{173D57}
\hypersetup{colorlinks=true,urlcolor=ink,linkcolor=ink,pdftitle={BGS-F. Количественные зазоры при фиксированном кике}}
\pagestyle{fancy}\fancyhf{}\lhead{\small BGS-F}\rhead{\small Количественные зазоры при фиксированном кике}\cfoot{\thepage}
\setlength{\headheight}{15pt}\setlength{\parindent}{1em}\setlength{\parskip}{2.5pt}
\setlength{\emergencystretch}{3em}\tolerance=3000
\newcommand{\Tr}{\operatorname{Tr}}
\newcommand{\EE}{\mathbb E}
\newcommand{\R}{\mathbb R}
\newcommand{\Z}{\mathbb Z}
\newcommand{\C}{\mathbb C}
\newcommand{\si}{\mathrm{Sine}_2}
\newcommand{\eps}{\varepsilon}
\newcommand{\sinc}{\operatorname{sinc}}
\newcommand{\sect}[1]{\section*{\sffamily\color{ink}#1}}
\newcommand{\subt}[1]{\par\medskip\noindent\textbf{#1.}\quad}
\newcommand{\T}{\mathbb T}
\newcommand{\dist}{\operatorname{dist}}
\newcommand{\Ad}{\operatorname{Ad}}
\newcommand{\zp}{\mathbb Z_7}
\newcommand{\QQ}{\mathbb Q}
\newcommand{\kk}{\overline{\mathbb Q}}
\newcommand{\ord}{\operatorname{ord}}
\newcommand{\diag}{\operatorname{diag}}
\begin{document}
\noindent{\small\sffamily\color{ink} МАТЕМАТИЧЕСКАЯ ФИЗИКА\hfill 22 НОЯБРЯ 2026}
\begin{center}{\LARGE\sffamily Количественные зазоры\\[3pt]при фиксированном кике}\\[6pt]
{\large Все нечётные размерности и точные интервалы при \(N=49\)}
\end{center}
Пусть \(U_N=\exp(iN10^{-80}A_N)M_N\) из исходной задачи BGS-F, а \(\lambda_1,\ldots,\lambda_r\) --- различные собственные числа полного оператора с кратностями \(d_1,\ldots,d_r\). Для любого нечётного \(N\ge3\) положим
\[
     R=N^3\log(10N),\qquad B_N=\exp\!\bigl(\exp(20R)\bigr).
\]

\subt{Эффективная граница зазора}
Если \(\lambda_i\ne\lambda_j\), то
\[
             |\lambda_i-\lambda_j|\ge \exp(-B_N).
\]
Оценка применима при неизвестных кратностях и к исходной силе без изменения параметра. Доказательство использует ненулевой определитель моментов ровно того размера, который равен числу различных собственных чисел. Различные алгебраические показатели приводятся к независимым с помощью целочисленной матрицы следового спаривания. Это даёт явную степень многочлена, к которому применяется эффективная теорема Линдемана--Вейерштрасса.

\subt{Равномерная оценка среднего обычного следа}
Для любых целых \(a\) и \(T\ge1\)
\[
 \left|\frac1{TN}\sum_{m=a+1}^{a+T}|\Tr U_N^m|^2
                   -\frac1N\sum_{j=1}^r d_j^2\right|
             \le\frac{2N}{T}\exp(B_N).
\]
Например, на \(N=7^k\) верхняя граница из теоремы о расщеплении сохраняется при совместном пределе, если \(T_N/(N\exp(B_N))\to\infty\). Прежняя формула о кратностях давала последовательный предел по времени и по размерности.

\subt{Сертифицированные зазоры полного оператора при \(N=49\)}
Для шести двойных узлов с параметрами \(\ell=4,6\) и \(h=0,1,2\) первый порядок скалярен. Получены строгие интервалы угловых зазоров при \(t_0=49/10^{80}\). Например, при \((\ell,h)=(4,0)\) зазор в радианах принадлежит
\[
          t_0^2[0.00170645,\,0.00170648].
\]
Интервалы вычисляются с направленным рациональным округлением. Остаток в дополнении Шура контролируется в операторной норме, поэтому интервалы относятся к собственным фазам полного \(U_{49}\).

Асимптотическая оценка зазоров лежит значительно ниже масштаба \(1/N\). Равномерность по началу временного окна относится к явно указанным большим \(T\). Предел Sine\(_2\) при \(10^{-80}\) и парный тест при временах порядка \(N\) остаются открытыми.

\newpage
\sect{1. Определитель моментов при неизвестной кратности}
Все обозначения \(A_N,M_N\) взяты из исходной постановки. Пусть \(B=4000A_N\) имеет различные собственные числа \(\eta_1,\ldots,\eta_s\) и спектральные проекторы \(P_1,\ldots,P_s\). Положим \(t=N/10^{80}\),
\[
 U(X)=\left(\sum_{j=1}^s X_j P_j\right)M_N,\qquad
 X^*_j=\exp(it\eta_j/4000).
\]
Тогда \(U(X^*)=U_N\). При \(r\) различных собственных числах \(U_N\) определим
\[
    H_r(X)=\bigl(\Tr U(X)^{a+b}\bigr)_{0\le a,b<r}.
\]
В точке \(X^*\) имеем факторизацию \(H_r=V\diag(d_j)V^{\mathsf T}\), \(V_{aj}=\lambda_j^a\). Поэтому
\begin{equation}
  \bigl|\det H_r(X^*)\bigr|
    =\prod_{j=1}^r d_j\prod_{1\le i<j\le r}|\lambda_i-\lambda_j|^2>0.
                                                        \tag{1}
\end{equation}
Индексы \(a,b\) начинаются с нуля, а степень каждого слагаемого определителя равна \(J=r(r-1)\). Формула (1) остаётся точной при любых кратностях. Значение \(r\) нужно только для доказательства, а численная граница будет от него независима.

\subt{Целый многочлен}
Пусть \(\Delta=\prod_{i<j}(\eta_i-\eta_j)^2\), где пустое произведение равно единице. Проектор выражается формулой
\[
    P_i=\prod_{j\ne i}\frac{B-\eta_jI}{\eta_i-\eta_j}.
\]
Поэтому \(\Delta P_i\) имеет целые алгебраические элементы. Матрица \(\sqrt N M_N\) тоже имеет целые алгебраические элементы. Значит
\[
      P(X)=(\sqrt N\Delta)^J\det H_r(X)
\]
имеет целые алгебраические коэффициенты. Множитель \(\sqrt N\Delta\) можно целиком перенести внутрь матрицы \(U(X)\) в каждой степени, что объясняет одинаковый знаменатель всех слагаемых.

\subt{Как один определитель ограничивает все пары}
Для любой выбранной пары \(i,j\) остальные хордовые расстояния в (1) не превосходят 2. Кроме того, \(\prod d_j\le N^N\). Поэтому при \(r\ge2\)
\[
  2\log|\lambda_i-\lambda_j|
          \ge \log|\det H_r(X^*)|-N\log N-N^2\log2.
                                                        \tag{2}
\]
Отдельная обработка неизвестных собственных чисел \(\lambda_i\) не требуется.

\newpage
\sect{2. Арифметика поля и высоты}
\subt{Общий контроль всех вложений}
Пусть \(K_0=\QQ(\zeta_{4N},\sqrt N)\), а \(K\) получается присоединением всех \(\eta_j\). Степень расширения
\[
   [K:\QQ]\le D:=8N N!
\]
следует из степени поля \(K_0\) и верхней границы \(N!\) для степени поля разложения. По исходным коэффициентам \(B\) сумма модулей не превосходит \(4744\). Каждое вложение \(K_0\) сохраняет унитарность \(M_N\) и эрмитовость \(B\). Действительно, сопряжение на \(\zeta_{4N}\) коммутирует с подстановкой любой другой первообразной степени, а образ \(\sqrt N\) вещественен. Каждая пара противоположных вейлевых меток остаётся парой сопряжённых матриц. Поэтому для всех вложений \(\sigma\) числа \(\sigma\eta_j\) вещественны и по модулю меньше \(H:=5000\). Проекторы \(\sigma(P_j)\) ортогональны, поскольку это спектральные проекторы \(\sigma(B)\).

\subt{Коэффициенты многочлена}
Во всех вложениях \(|\sigma\Delta|\le(2H)^{N(N-1)}\). Сумма модулей коэффициентов в \(\sigma(P)\) ограничена сверху
\[
  r!N^r\bigl(s\sqrt N\,|\sigma\Delta|\bigr)^J.
\]
Для этого матричный многочлен \((\sqrt N\Delta)U(X)\) оценивается по сумме операторных норм коэффициентов. Эта сумма умножается при произведении, применение следа стоит множитель \(N\), а в определителе не больше \(r!\) слагаемых. Указанная верхняя граница не превосходит \(\exp(b_\beta)\), где
\[
       b_\beta=2N\log N+
                N^2\bigl(2\log N+N^2\log(2H)\bigr).
\]
Та же оценка справедлива для логарифма абсолютного значения \((\sqrt N\Delta)^J\). Все коэффициенты целые, поэтому \(b_\beta\) ограничивает совместную логарифмическую высоту в любом из полей, содержащих коэффициенты.

\subt{Дискриминант подполя}
В \(K\) выберем базис над \(\QQ\) из целых мономов в \(\zeta_{4N},\sqrt N,\eta_1,\ldots,\eta_s\). Достаточно степеней \(0,1\) для \(\sqrt N\) и степеней от \(0\) до \(N-1\) для каждой \(\eta_j\), поскольку последняя удовлетворяет монному характеристическому многочлену над \(K_0\). Модуль каждого сопряжения базисного монома не больше \(\sqrt N\,H^{N^2}\). По неравенству Адамара для квадрата определителя вложений
\[
  \frac{\log|\Delta_K|}{[K:\QQ]}
                 \le \log D+\log N+2N^2\log H=:b_\Delta .
\]
Дискриминант целого базиса поля делит дискриминант выбранного набора целых элементов, поскольку последний отличается квадратом целого индекса. Для подполя \(F\subset K\) та же граница на \(\log|\Delta_F|/[F:\QQ]\) следует из формулы дискриминантов в башне. Она охватывает поле всех коэффициентов после замены переменных.

\newpage
\sect{3. Независимые экспоненты и эффективная оценка}
Выберем базис \(b_1,\ldots,b_m\) пространства над \(\QQ\), порождённого \(\eta_1,\ldots,\eta_s\), из самих \(\eta_j\). Пусть \(F_0=\QQ(\eta_1,\ldots,\eta_s)\). Все его вложения вещественны по разделу 2. Поэтому целочисленная матрица
\[
   C_{ij}=\Tr_{F_0/\QQ}(b_i b_j)
\]
положительно определена при \(m>0\). Для ненулевой рациональной комбинации квадрат имеет положительный след, а затем положительность следует из невырожденности спаривания. Положим \(Q=\det C>0\). По правилу Крамера \(\eta_j=Q^{-1}\sum_i n_{ij}b_i\), где \(Q,n_{ij}\) целые и по модулю не превосходят
\[
               W=(N D H^2)^N.
\]
Действительно, каждый след попарного произведения по модулю не больше \(D H^2\), и та же оценка относится к столбцу правых частей. В терминах
\[
  \alpha_i=itb_i/(4000Q)
\]
замена \(X_j=\prod_iY_i^{n_{ij}}\), \(Y_i=e^{\alpha_i}\), даёт точные значения \(X_j^*\). В каждом показателе после подстановки абсолютное значение не больше \(N^2 W\). Домножаем на \(\prod_iY_i^{N^2W}\) и получаем целый многочлен \(\widetilde P\) степени \(\le L=2N^3 W\). При вычислении модуль домноженного монома равен единице. Сумма модулей коэффициентов не увеличивается. Поскольку \(P(X^*)\ne0\), новый многочлен ненулевой. Числа \(\alpha_i\) независимы над \(\QQ\). Их совместная высота ограничена
\[
 b_\alpha=\log(4000\cdot10^{80}WNH).
\]
Общий знаменатель не больше \(4000\cdot10^{80}W\), числители целые, а все архимедовы модули чисел \(\alpha_i\) не превосходят \(NH\).

\subt{Применение теоремы Серта}
Используем эффективную теорему Линдемана--Вейерштрасса в многочленной форме, приведённой как теорема 5.1 у Хуана~[1]. При ограничениях выше её явные константы \(r,r',r''\) сверху ограничены числом
\[
                  C_0=10^4(9ND)^{N+2}(1+b_\alpha+\log(9ND)).
\]
Для проверки подставляем \(m\le N\), \([K:\QQ]\le D\), \(h(\alpha)\le b_\alpha\) во все три константы. Общий степенной множитель \(3^{2m}m^m[K:\QQ]^{m+1}\) не превосходит \((9ND)^ND\). Множители \(1+6[K:\QQ]\) после выделения одного сомножителя \([K:\QQ]\) дают не более 7. Множитель \(1+6m\) поглощается дополнительными степенями \((9ND)^2\). Логарифмы и высота учтены в последней скобке. С поправкой на дискриминант и оценкой \(\max|\alpha_i|\le NH\) получается
\begin{equation}
 -\log|P(X^*)|\le C_0 L^N
 \left(b_\beta+b_\Delta+\exp\bigl(C_0 L^N(1+\log L)+72NH\bigr)\right).
                                                        \tag{3}
\end{equation}
Если \(\widetilde P\) постоянен, нижнюю границу можно также получить непосредственно из нормы ненулевого целого коэффициента. При \(m=0\) имеем \(A_N=0\). Тогда все отношения фаз являются корнями единицы порядка, делящего порядок \(S\bmod N\), не больше \(N^4\). Такой случай удовлетворяет заявленной границе сразу.

\newpage
\sect{4. Явная скорость временного усреднения}
Покажем упрощение (3). При \(u=\log(10N)\) и \(R=N^3u\) получаем
\[
 \log D\le2Nu,\quad
 \log W\le5N^2u,\quad
 \log L\le6N^2u,\quad
 b_\alpha\le13N^2u,\quad
 \log C_0\le8N^2u .
\]
Первые три оценки используют \(N!\le N^N\), \(H^2\le (10N)^6\) и \(N\ge3\). Далее \(\log(4000\cdot10^{80})\le7N^2u\), \(\log(NH)\le4u\) и \(1+b_\alpha+\log(9ND)\le16N^2u\). Подстановка даёт последнюю оценку. Кроме того, \(\log(1+\log L)\le4u\). Таким образом, \(C_0L^N\le e^{9R}\), а \(C_0 L^N(1+\log L)+72NH\le e^{11R}\). Оставшиеся полиномиальные слагаемые ограничены сверху \(e^R\):
\[
       b_\beta+b_\Delta\le e^R,\qquad
       b_\beta+N\log N+N^2\log2\le e^R .
\]
Например, обе левые части не превосходят \(20N^4u\), что меньше \(e^R\) при \(N\ge3\). Отсюда правая часть (3) не больше \(\exp(\exp(12R))\). Чтобы перейти к \(|\det H_r|\), вычитаем логарифм масштабного множителя, ограниченный \(b_\beta\). Затем применяем (2). Получаем \(|\lambda_i-\lambda_j|\ge e^{-B_N}\) с единым \(B_N=\exp(\exp(20R))\). При \(r=1\) утверждение о зазорах пусто, а формула усреднения точна при каждом \(T\).

\subt{Равномерное усреднение}
Тождество
\[
      |\Tr U_N^m|^2=\sum_{i,j=1}^r d_i d_j(\lambda_i/\lambda_j)^m
\]
даёт диагональный вклад \(\sum d_j^2\) при любом \(m\). Для \(i\ne j\) модуль геометрической суммы по \(a+1,\ldots,a+T\) не превосходит \(2/|1-\lambda_i/\lambda_j|\), независимо от \(a\). Поэтому остаток в среднем ограничивается числом
\[
      \frac{2e^{B_N}}{TN}\sum_{i\ne j}d_i d_j
                    \le\frac{2Ne^{B_N}}T .
\]
На \(N=7^k\) уже доказано \(\limsup N^{-1}\sum d_j^2\le58663/57624\). Теперь если \(T_N/(N e^{B_N})\to\infty\), в этом верхнем пределе можно заменить \(N^{-1}\sum d_j^2\) указанным средним по любому начальному \(a_N\).

\subt{Масштаб оценки}
Хордовая граница даёт угловую границу на окружности: для фаз \(\theta_i\in\R/\Z\) имеем \(\operatorname{dist}(\theta_i-\theta_j,\Z)\ge e^{-B_N}/(2\pi)\). Поэтому внутри меньшего интервала могут встретиться только точные совпадения. Интервал длины \(1/N\) допускает множество различных фаз и требует другой оценки.

\newpage
\sect{5. Точный малый зазор на исходной силе}
Возьмём \(N=49\), \(m=21\), \(t=49/10^{80}\), \(A=A_{49}\). После унитарной симплектической замены \(M\) становится дилатацией \(\psi(x)\mapsto\psi(\gamma^{-1}x)\) с точностью до общей фазы. Пусть \(\lambda\) --- двойной линейный уровень при \(\ell=4\) или \(6\) и \(h=0,1,2\), \(E\) его проектор. Во всех шести узлах \(EAE=cE\) точно по целому циклотомическому расчёту. Пишем \(Q=I-E\) и вводим эрмитову матрицу на \(E\)
\[
    H=\frac{1}{2i}EAQ\frac{\lambda+M}{\lambda-M}QAE .
\]
Норма \(A\) меньше 2. Расстояние от \(\lambda\) до остальных узлов \(g\ge4/m\). Следовательно, \(\|H\|\le4/g\le21\). Подставим
\[
      z=\lambda e^{itc}(1+it^2x),\qquad |x|\le22 .
\]
Пусть \(S(z,t)\) --- дополнение Шура матрицы \(U_{49}-zI\) на \(E\). На указанной области
\[
        \frac{S(z,t)}{-i\lambda t^2}
                 =xI+H+\mathcal E(x,t),\qquad
                     \|\mathcal E(x,t)\|\le10^4t.             \tag{4}
\]

\subt{Проверка остатка}
Действительно, \(|z-\lambda|\le3t\) и \(\|U_{49}-M\|\le2t\). Обратный дальний блок имеет норму не больше \(2/g\), а его отличие от обратной матрицы для \(M-\lambda I\) не больше \(10t/g^2\). Для каждого внедиагонального блока ошибка первого порядка не больше \(2t^2\), а сама норма блока не больше \(2t\). Поэтому в смешанном произведении после вычитания второй степени остаток ограничен числом \((40/g^2+16/g)t^3\). В диагональном блоке и в \(z\) оставшиеся тейлоровы члены в сумме не больше \(60t^3\). После замены обратной матрицы коэффициент при \(-\lambda t^2\) равен \(EA^2E/2+EAQM(\lambda-M)^{-1}QAE=c^2 E/2+iH\). При вычитании \(zE\) общий член с \(c^2E/2\) сокращается. Это доказывает (4).

\subt{Переход от коэффициента к истинному спектру}
Пусть \(h_1<h_2\) --- собственные числа \(H\). Рациональный расчёт на следующей странице показывает \(h_2-h_1>0.001\). На каждой из окружностей радиуса \(10^{-10}\) около \(-h_j\) обратная матрица \((xI+H)^{-1}\) имеет норму не больше \(10^{10}\), а \(10^4t<10^{-10}\). Значит, при непрерывном включении остатка в (4) определитель сохраняет число нулей внутри каждой окружности. Оно равно одному. Так получаются две фазы полного оператора. Их разность в радианах, делённая на \(t^2\), отличается от \(h_2-h_1\) менее чем на \(10^{-8}\). Действительно, радиусы окружностей дают погрешность не больше \(2\cdot10^{-10}\) для разности вещественных частей \(x\), а при \(|x|\le22\) ошибка \(t^{-2}\arg(1+it^2x)-\Re x\) по модулю не больше \(2\cdot22^2t^2\).

\newpage
\sect{6. Целочисленные интервалы и воспроизведение}
Для вычисления \(H\) удобно разложить рациональную функцию от дилатации. Если \(Z=M/\lambda\), то \(Z^m=I\), и
\[
    Q\frac{1+Z}{1-Z}Q
                  =\sum_{b=1}^{m-1}(1-2b/m) Z^b .
\]
На \(Z=1\) правая часть равна нулю. На остальных корнях формула следует дифференцированием конечной геометрической суммы. Для сжатия используем два ненормированных орбитальных вектора, деля скалярные произведения на \(m\). Фазы \(M^b/\lambda^b\) и элементы \(A\) вычисляются непосредственно по исходным вейлевым меткам. Полученная матрица имеет размер \(2\times2\). Её зазор равен
\[
       \Delta_H=\sqrt{(H_{11}-H_{22})^2+4|H_{12}|^2}.
\]
Все квадраты и скалярные произведения посчитаны в рациональных интервалах. На сетке \(10^{-50}\) операции округляются наружу. Границы \(\pi\) получаются формулой Мэчина и знакопеременным рядом для арктангенса. Для косинуса и синуса остаток ряда Тейлора ограничивается сверху \(4^{89}/89!<10^{-50}\), поскольку аргумент приведён к \([-\pi,\pi]\). Для квадратного корня используются целочисленные нижняя и верхняя границы.

\begin{center}
\begin{tabular}{ccll}
\toprule
\(\ell\)&\(h\)& Интервал для \(\Delta_H\)& Интервал зазора в радианах, делённого на \(t^2\)\\
\midrule
4&0&\([0.00170646,0.00170647]\)&\([0.00170645,0.00170648]\)\\
4&1&\([0.00276813,0.00276814]\)&\([0.00276812,0.00276815]\)\\
4&2&\([0.00191625,0.00191626]\)&\([0.00191624,0.00191627]\)\\
6&0&\([0.01244094,0.01244095]\)&\([0.01244093,0.01244096]\)\\
6&1&\([0.00136920,0.00136921]\)&\([0.00136919,0.00136922]\)\\
6&2&\([0.01286096,0.01286097]\)&\([0.01286095,0.01286098]\)\\
\bottomrule
\end{tabular}
\end{center}

\subt{Сверка}
Команда \texttt{python3 effective\_gaps/check\_intervals.py --verify} пересчитывает и побайтно проверяет \texttt{effective\_gaps/data/intervals.json}. Она проверяет условие первого порядка на шести узлах и использует только стандартную библиотеку Python. Скрипт \texttt{effective\_gaps/check\_matrix.py} независимо вычисляет вторую матрицу через плотное обращение дальнего блока. Он сверяет все шесть численных коэффициентов с рациональными интервалами. Вывод об интервалах истинных фаз опирается на целочисленный расчёт и (4).

\subt{Источник эффективной оценки}
[1] C.-C. Huang. \emph{Explicit Bounds for Linear Forms in the Exponentials of Algebraic Numbers}, ISSAC 2022, теорема 5.1, многочленная форма теоремы A. Sert. \href{https://arxiv.org/abs/2112.05004}{arXiv:2112.05004}.

\end{document}
